\title{Megalodon: Efficient LLM Pretraining and Inference with Unlimited Context Length}

\begin{abstract}
The quadratic complexity and weak length extrapolation of Transformers limits their ability to scale to long sequences, and while sub-quadratic solutions like linear attention and state space models exist, they empirically underperform Transformers in pretraining efficiency and downstream task accuracy. We introduce \textsc{Megalodon}, an neural architecture for efficient sequence modeling with unlimited context length.
\textsc{Megalodon} inherits the architecture of \textsc{Mega} (exponential moving average with gated attention), and further introduces multiple technical components to improve its capability and stability, including \emph{complex exponential moving average (CEMA)}, \emph{timestep normalization} layer, \emph{normalized attention} mechanism and \emph{pre-norm with two-hop residual} configuration. In a controlled head-to-head comparison with \textsc{Llama}2, 
\textsc{Megalodon} achieves better efficiency than Transformer in the scale of 7 billion parameters and 2 trillion training tokens.
\textsc{Megalodon} reaches a training loss of 1.70, landing mid-way between \textsc{Llama}2-7B (1.75) and 13B (1.67). 

The improvements of \textsc{Megalodon} over Transformers are robust throughout a range of benchmarks across different tasks and modalities. \\
\textbf{Code}: \url{https://github.com/XuezheMax/megalodon}
\end{abstract}

\section{Introduction}
In many real-world applications, such as multi-turn conversation, long-document comprehension, and video generation, large language models (LLMs) must efficiently process long sequential data, understand internal long-range dynamics, and generate coherent output. The Transformer architecture~\citep{Vaswani+2017}, despite its remarkable capabilities, faces challenges with quadratic computational complexity and limited inductive bias for length generalization, making it inefficient for long sequence modeling~\citep{wang2024limits,zhou2024transformers}. Even with recently proposed distributed attention solutions~\citep{li2023lightseq,liu2024ring}, computing a single training step of a 7B parameter model over a 1M-token sequence is more than 100 times slower than performing the equivalent computation using 256 separate sequences of 4K tokens each.

Techniques like efficient attention mechanisms~\citep{tay2020efficient,ma2021luna} and structured state space models~\citep{gu2022efficiently,poli2023hyena,gu2023mamba} have been introduced to overcome these limitations, aiming to enhance scalability and performance. However, the practical application of these methods still falls short of Transformers~\citep{tay2022scaling,gu2023mamba}. 

This work introduces an unlimited context model that outperforms the canonical Transformer architecture on real-world language modeling.

We introduce \textsc{Megalodon}, an improved \textsc{Mega} architecture~\citep{ma2023mega}, which harnesses the gated attention mechanism with the classical exponential moving average (EMA)~\citep{hunter1986exponentially} approach (\S\ref{sec:background}).
To further improve the capability and efficiency of \textsc{Megalodon} on large-scale long-context pretraining, we propose multiple novel technical components. First, \textsc{Megalodon} introduces the \emph{complex exponential moving average (CEMA)} component, which extends the multi-dimensional damped EMA in \textsc{Mega} to the complex domain (\S\ref{subsec:cema}).
Then, \textsc{Megalodon} proposes the \emph{timestep normalization} layer, which generalizes the group normalization layer~\citep{wu2018group} to auto-regressive sequence modeling tasks to allow normalization along the sequential dimension (\S\ref{subsec:timestepnorm}).
To improve large-scale pretraining stability, \textsc{Megalodon} further proposes \emph{normalized attention} (\S\ref{subsec:norm-attn}), together with \emph{pre-norm with two-hop residual} configuration by modifying the widely-adopted pre- and post-normalization methods (\S\ref{subsec:pre-post-norm}).
By simply chunking input sequences into fixed blocks, as is done in  \textsc{Mega}-chunk~\citep{ma2023mega}, \textsc{Megalodon} achieves linear computational and memory complexity in both model training and inference.

Empirically, we demonstrate the potential of \textsc{Megalodon} as a general architecture for modeling long sequences, by evaluating its performance across multiple scales of language modeling, as well as downstream domain-specific tasks. Through a direct comparison with \textsc{Llama2}, while controlling for data and compute, \textsc{Megalodon}-7B significantly outperforms the state-of-the-art variant of Transformer used to train \textsc{Llama2}-7B~\citep{touvron2023llama} on both training perplexity (Figure~\ref{fig:train_loss}) and across downstream benchmarks (Table~\ref{tab:benmarks}).
Evaluation on long-context modeling, including perplexity in various context lengths up to 2M and long-context QA tasks in Scrolls~\citep{parisotto2020stabilizing} prove \textsc{Megalodon}'s ability to model sequences of unlimited length. Additional experimental results on small/medium-scale benchmarks, including LRA~\citep{tay2021long}, ImageNet~\citep{deng2009imagenet}, Speech Commands~\citep{warden2018speech}, WikiText-103~\citep{merity2017pointer} and PG19~\citep{raecompressive2019}, demonstrate the robust improvements of \textsc{Megalodon} across scales and modalities.

\section{Background: Moving Average Equipped Gated Attention (\textsc{Mega})}
\label{sec:background}
In this section, we setup notations, briefly review the key components in the \textsc{Mega} architecture~\citep{ma2023mega}, and discuss the existing problems in \textsc{Mega}.

Following the notations in \textsc{Mega}, we use $\boldsymbol{X} = \{\mathbf{x}_1, \mathbf{x}_2, \ldots, \mathbf{x}_n\} \in \mathbb{R}^{n\times d}$ and $\boldsymbol{Y} =\{\mathbf{y}_1, \mathbf{y}_2, \ldots, \mathbf{y}_n\} \in \mathbb{R}^{n\times d}$ to denote the input and output sequences with length $n$, and assume the representations of the input and output sequences have the same dimension $d$. 

\subsection{Multi-dimensional Damped EMA}
\textsc{Mega} embeds an EMA component into the calculation of the attention matrix to incorporate inductive biases across the timestep dimension. Concretely, the multi-dimensional damped EMA first expands each dimension of the input sequence $\boldsymbol{X}$ individually into $h$ dimensions via an expansion matrix $\boldsymbol{\beta} \in \mathbb{R}^{d\times h}$, then applies damped EMA to the $h$-dimensional hidden space. Formally, for each dimension $j \in \{1, 2, \ldots, d\}$:

\begin{align}
\mathbf{u}^{(j)}_t & = \boldsymbol{\beta}_j \mathbf{x}_{t,j} \nonumber \\
\mathbf{h}^{(j)}_t & = \boldsymbol{\alpha}_j \odot \mathbf{u}^{(j)}_t + (1 - \boldsymbol{\alpha}_j \odot \boldsymbol{\delta}_j) \odot \mathbf{h}^{(j)}_{t-1} \label{eq:mddema} \\
\mathbf{y}_{t,j} & = \boldsymbol{\eta}^T_j \mathbf{h}^{(j)}_t \nonumber
\end{align}

where $\mathbf{u}^{(j)}_t \in \mathbb{R}^{h}$ is the expanded $h$-dimensional vector for the $j$-th dimension at timestep $t$. 
$\boldsymbol{\alpha} \in (0, 1)^{d\times h}$, $\boldsymbol{\delta} \in (0, 1)^{d\times h}$ are the decaying and damping factors, respectively. $\mathbf{h}^{(j)}_t \in \mathbb{R}^{h}$ is the EMA hidden state for the $j$-th dimension at timestep $t$. 
$\boldsymbol{\eta} \in \mathbb{R}^{d\times h}$ is the projection matrix to map the $h$-dimensional hidden state back to $1$-dimensional output $\mathbf{y}_{t,j} \in \mathbb{R} $.

\subsection{Moving Average Equipped Gated Attention} In the gated attention mechanism in \textsc{Mega}, the output from EMA \eqref{eq:mddema} is used to compute the shared representation~\citep{hua2022transformer}, because it encodes contextual information through EMA. Subsequently, \textsc{Mega} introduces the reset gate, the update gate , and computes the candidate activation with the update gate and the residual connection. The technical details are provided in Appendix~\ref{appendix:mega}.

\subsection{Existing Problems in \textsc{Mega}}
To reduce the quadratic complexity in the full attention mechanism, \textsc{Mega} simply split the sequences of queries, keys and values in (\ref{eq:query}-\ref{eq:value}) into chunks of length $c$. The attention in \eqref{eq:attention2} is individually applied to each chunk, yielding linear complexity $O(kc^2)=O(nc)$. Technically, the EMA sub-layer in \textsc{Mega} helps capture local contextual information near each token, mitigating the problem of losing contextual information beyond chunk boundaries in the chunk-wise attention.

Despite the impressive successes of \textsc{Mega}, it still suffers its own problems: i) the performance of \textsc{Mega} with chunk-wise attention still fails behind the one with full attention, due to the limited expressiveness of the EMA sub-layer in \textsc{Mega}. 
ii) for different tasks and/or data types, there are architectural divergences in the final \textsc{Mega} architectures. For example, different normalization layers, normalization patterns~(pre-norm vs. post-norm) and attention functions ($f(\cdot)$ in \eqref{eq:attention2}) are applied to different data types~(see \citet{ma2023mega} for details). iii) There are no empirical evidences showing that \textsc{Mega} is scalable for large-scale pretraining.

\section{\textsc{Megalodon}}
\label{sec:mega2}
To address the aforementioned problems of \textsc{Mega}, in this section we describe the novel technical advancements of \textsc{Megalodon}.

\subsection{CEMA: Extending Multi-dimensional Damped EMA to Complex Domain}
\label{subsec:cema}
As discussed in \citet{ma2023mega}, the EMA component can be regarded as a simplified state space model with diagonal state matrix. Directly inspired from \citet{gu2022parameterization}, as almost all matrices diagonalize over the complex plane, a straight-forward idea to improve EMA capability is to extend to work over the complex number system $\mathbb{C}$. 
We propose the \emph{complex exponential moving average (CEMA)}, which re-writes Eq.~\eqref{eq:mddema}:

\begin{align}
\label{eq:cema}
\mathbf{h}^{(j)}_t & = \boldsymbol{\alpha}_j (\cos{\theta_j + i \sin{\theta_j}}) \odot \mathbf{u}^{(j)}_t + (1 - \boldsymbol{\alpha}_j \odot \boldsymbol{\delta}_j) (\cos{\theta_j + i \sin{\theta_j}}) \odot \mathbf{h}^{(j)}_{t-1} \nonumber \\
\mathbf{y}_{t,j} & = \mathrm{Re}(\boldsymbol{\eta}^T_j \mathbf{h}^{(j)}_t )
\end{align}

where $\boldsymbol{\alpha}$, $\boldsymbol{\delta} \in \mathbb{R}^{d\times h}$ are the real number parameters same as in EMA. Different from EMA, $\boldsymbol{\eta} \in \mathbb{C}^{d\times h}$ in CEMA are complex numbers. 
$\theta_j \in \mathbb{R}^{h}, \,\, j \in \{1, 2, \ldots, d\}$ are the $h$ arguments. To uniformly space the $h$ arguments over the period $2\pi$, we parameterize $\theta_j$ as:

\begin{equation}
\theta_{j,k} = \frac{2\pi k}{h} \omega_j, \quad \forall k \in \{1, 2, \ldots, h\}
\end{equation}

where the learnable parameter $\omega \in \mathbb{R}^{d}$ depicts the $d$ base angles. By decaying the absolute value of each $h_t$, CEMA preserves the decaying structure in kernel weights, which is a key principle to the success of convolutional models on long sequence modeling~\citep{li2023makes}.

\subsection{Timestep Normalization}
\label{subsec:timestepnorm}

Despite the impressive performance of Layer Normalization combined with Transformer, it is obvious that layer normalization cannot directly reduce the internal covariate shift along the spatial dimension (a.k.a timestep or sequential dimension)~\citep{ioffe2015batch}. Group Normalization~\citep{wu2018group} normalizes hidden representations both along the timestep dimension and a subset of the feature dimension, which has obtained improvements over Layer Normalization on a range of computer vision tasks. However, it cannot be directly applied to Transformer on auto-regressive sequence modeling, due to the leakage of future information via the mean and variance across the timestep dimension. 

In \textsc{Megalodon}, we extend Group Normalization to the auto-regressive case by computing the cumulative mean and variance. Formally, suppose an input sequence $\boldsymbol{X} = \{\mathbf{x}_1, \mathbf{x}_2, \ldots, \mathbf{x}_n\} \in \mathbb{R}^{n\times d}$, and $k$ groups along the feature dimension with $d_g = d/k$ elements per group. Then, the mean and variance of the first group at timestep $t \in \{1, 2, \ldots, n \}$ are:

\begin{equation}
\mu_t = \frac{1}{t * d_g}\sum\limits_{i = 1}^{t}\sum\limits_{j=1}^{d_g} x_{i,j}, \qquad \sigma^2_t = \frac{1}{t * d_g}\sum\limits_{i = 1}^{t}\sum\limits_{j=1}^{d_g} (x_{i,j} - \mu_t)^2
\end{equation}

Figure~\ref{fig:normalization} illustrates Layer Normalization and Timestep Normalization. To efficiently and precisely calculate the cumulative mean and variance in each timestep, we provide hardware-friendly implementation on modern hardware (GPU) (see Appendix~\ref{appendix:cuda}). 

\subsection{Normalized Attention in \textsc{Megalodon}}
\label{subsec:norm-attn}

Previous studies have investigated the saturation and instability issues in the original scaled dot-product attention \eqref{eq:attention2}. A number of novel techniques have emerged to modify the scaled dot-product attention, among which normalized attention mechanisms, such as (scaled-) cosine attention~\citep{luo2018cosine,liu2022swin} and QK-normalization~\citep{henry2020query}, have stood out for the simplicity and effectiveness. 

Directly inspired from these normalized attention mechanisms, we propose the normalized attention mechanism specifically defined for \textsc{Mega} to improve its stability. Formally,

\begin{align}
\boldsymbol{X}' & = \mathrm{CEMA}(\boldsymbol{X}) \qquad & \qquad \in \mathbb{R}^{n\times d} \\
\boldsymbol{Z} & = \boldsymbol{X}' W_z + b_z, \quad \boldsymbol{Z}' = \frac{\boldsymbol{Z}}{\| \boldsymbol{Z}\|} & \qquad \in \mathbb{R}^{n\times z} \\
\boldsymbol{Q} & = \boldsymbol{\kappa}_q \odot \boldsymbol{Z}' + \boldsymbol{\mu}_q \qquad & \qquad \in \mathbb{R}^{n\times z} \label{eq:q} \\
\boldsymbol{K} & = \boldsymbol{\kappa}_k \odot \boldsymbol{Z}' + \boldsymbol{\mu}_k \qquad & \qquad \in \mathbb{R}^{n\times z} \label{eq:k}
\end{align}

where $\boldsymbol{Q}$ and $\boldsymbol{K}$ are computed by using the normalized shared representation $\boldsymbol{Z}'$ instead of $\boldsymbol{Z}$. Note that we remove the SiLU~\citep{ramachandran2017swish} activation function $\phi_{\mathrm{silu}}$ in \eqref{eq:z}, because the normalization on $\boldsymbol{Z}$ has incorporated non-linearity into $\boldsymbol{Z}'$.
Then the attention operation in \eqref{eq:attention2} has been changed to:

\begin{align}\label{eq:norm-attention}
\boldsymbol{O} & = f_{\mathrm{softmax}} \left(\boldsymbol{Q}{\boldsymbol{K}}^{T} \right) \boldsymbol{V} \quad & \qquad \qquad \in \mathbb{R}^{n\times v}
\end{align}

As we use learnable $\boldsymbol{\kappa}_q$, $\boldsymbol{\kappa}_k$ in \eqref{eq:q} and \eqref{eq:k}, we can remove the scaled term $\tau(\boldsymbol{X})$.
In addition, we found that with the normalized attention, the softmax function $f_{\mathrm{softmax}}$ obtains the best or at least comparable performance on different tasks and data modalities~(see Appendix~\ref{appendix:experiments}).
Hence, throughout this paper we use softmax as the default attention function.

\subsection{Pre-Norm with Two-hop Residual}
\label{subsec:pre-post-norm}
Normalization configurations are crucial in stably training deep architectures, and pre-normalization~\citep{xiong2020layer} has become the default normalization configuration because of its better convergence properties than post-normalization in the original Transformer architecture~\citep{Vaswani+2017}. However, extensive studies have investigated the instability issue of pre-normalization when scaling up model size~\citep{davis2021catformer,liu2022swin}. Formally, a Transformer-based block in pre-noromalization can be formulated as (shown in Figure~\ref{fig:mega2} (b)):

\begin{align}
\hat{\boldsymbol{Y}} & = \mathrm{Attention}(\mathrm{Norm}(\boldsymbol{X})) + \boldsymbol{X} \nonumber \\
\boldsymbol{Y} & = \mathrm{FFN}(\mathrm{Norm}(\hat{\boldsymbol{Y}})) + \hat{\boldsymbol{Y}} \nonumber \\
 & = \mathrm{FFN}(\mathrm{Norm}(\hat{\boldsymbol{Y}})) + \mathrm{Attention}(\mathrm{Norm}(\boldsymbol{X})) + \boldsymbol{X}
\end{align}

where the output $\boldsymbol{Y}$ is the sum of the input $\boldsymbol{X}$ and the output of each component in one block. Hence, the range and/or variance of $\boldsymbol{Y}$ keeps increasing for deeper blocks, causing the instability issue. In the original \textsc{Mega} architecture, the update gate $\boldsymbol{\varphi}$ \eqref{eq:update} is used for a gated residual connection \eqref{eq:output} to mitigate this problem~\citep{parisotto2020stabilizing,xu2020transformer}. However, the update gate $\boldsymbol{\varphi}$ introduces more model parameters and the instability issue still exists when scaling up model size to 7 billion.

\textsc{Megalodon} introduces a new configuration named \emph{pre-norm with two-hop residual}, which simply re-arranges the residual connections in each block (shown in Figure~\ref{fig:mega2} (c):

\begin{align}
\hat{\boldsymbol{Y}} & = \mathrm{Attention}(\mathrm{Norm}(\boldsymbol{X})) + \boldsymbol{X} \nonumber \\
\boldsymbol{Y} & = \mathrm{FFN}(\mathrm{Norm}(\hat{\boldsymbol{Y}})) + \boldsymbol{X}
\end{align}

where the input $\boldsymbol{X}$ is reused as the residual connection of the FFN layer. Since $\hat{\boldsymbol{Y}}$ is directly followed by a normalization layer, we remove the update gate $\boldsymbol{\varphi}$ and use standard residual connection. The graphical architecture of a \textsc{Megalodon} sub-layer is visualized in Figure~\ref{fig:mega2} (a). Note that the Timestep Normalization is only applied before the attention layer. Before the FFN layer, we still use Layer Normalization. The reasons are two-fold: i) Layer Normalization is faster than Timestep Normalization; ii) the output vector of each token from the attention layer is a mixture of vectors from contextual tokens via attention weights. Hence, normalizing the attention output along the feature dimension is similar to indirectly normalize along the timestep dimension. 

\subsection{4-Dimensional Parallelism in Distributed LLM Pretraining}
\label{subsec:parallelism}
Efficient distributed training algorithm is essential to train a large-scale language model, and several parallelization mechanisms have been introduced. The three most commonly used parallelism strategies are data, tensor~\citep{shoeybi2019megatron} and pipeline parallelism~\citep{NEURIPS2019_093f65e0}. However, the 3-dimensional parallelism is still insufficient to scale up the context length of LLMs~\citep{li2023lightseq,liu2024ring}.

Benefiting from the chunk-wise attention in \textsc{Megalodon}, we can efficiently parallelize it along the new timestep/sequence dimension, which is orthogonal to all the aforementioned three parallelism dimensions. In \textsc{Megalodon}, the only communications between devices in one chunk-parallel group are the last hidden state of CEMA and the cumulative mean and variance of Timestep Normalization in each block. Using asynchronous communication, we can minimize the overhead of chunk parallelization by hiding the communication costs in the computation of other components inside the same block and/or other blocks.

\section{Experiments}
\label{sec:exp}

To evaluate the scalability and efficiency of \textsc{Megalodon} on long-context sequence modeling, we scale up \textsc{Megalodon} to 7-billion model size and apply it to large-scale language model pretraining on 2 trillion tokens. We also conduct experiments on small/medium-scale sequence modeling benchmarks, including Long Range Arena (LRA)~\citep{tay2021long}, raw speech classification on Speech Commands~\citep{warden2018speech}, image classification on ImageNet-1K~\citep{deng2009imagenet}, and language-modeling on WikiText-103~\citep{merity2017pointer} and PG19~\citep{raecompressive2019}.~\footnote{Some results are provided in Appendix~\ref{appendix:experiments}, due to space limits.} Empirically, \textsc{Megalodon} significantly outperforms all the state-of-the-art baseline models on these tasks across various data modalities. 

\subsection{LLM Pretraining}
\paragraph{Architectural Details} In our \textsc{Megalodon}-7B model, we adopt most of architectural hyperparameters from \textsc{Llama}2-7B to ensure fair comparison: \textsc{Megalodon}-7B consists of 32 blocks, with feature dimension $d=4096$. Following \textsc{Llama}2, we use the SwiGLU activation function~\citep{shazeer2020glu} in the feed-forward layer, and rotary positional embedding~(RoPE, \citet{su2021roformer}).
We set the attention chunk size $c=4096$, which is the same as the pretraining context length in \textsc{Llama}2.
Benefiting from the attention gate ($\gamma$ in \eqref{eq:reset}), we use a much smaller number of attention heads $h=4$ in \textsc{Megalodon}-7B, comparing to $h=32$ in \textsc{Llama}2-7B. 
In addition, we apply pre-norm with two-hop residual (\S\ref{subsec:pre-post-norm}), using Timestep Normalization (\S\ref{subsec:timestepnorm}) and Layer Normalization~\citep{ba2016layer}, while \textsc{Llama}2 models apply pre-normalization with RMSNorm~\citep{zhang2019root}.

\paragraph{Data and Pretraining Details} We use the same mix of publicly available data from \textsc{Llama}2, ensuring that the model are trained on exactly the same 2-trillion tokens. We also use the same tokenizer as \textsc{Llama}2, whose vocabulary size is $32$K.

We trained \textsc{Megalodon}-7B using the AdamW optimizer~\citep{loshchilov2019decoupled}, with $\beta_1=0.9$, $\beta_2=0.95$, $\epsilon=1e-8$. The learning rate is $3.5e-4$ and cosine learning rate schedule is applied with warmup of $2500$ steps. We use a weight decay of $0.1$ and gradient clipping of $1.0$, and no dropout is applied during training. The context length in pretraining is $32$K (4 attention chunks). The global batch size is 4M tokens, and is distributed on 256 NVIDIA A100 GPUs (16K tokens per A100). We set data parallel size to 128, chunk parallel size to 2 and tensor parallel size to 1.

\paragraph{Data and Computation Efficiency} We evaluate the efficiency of \textsc{Megalodon} w.r.t both the data and computation perspectives. For data efficiency, we display the negative log-likelihood (NLL) for \textsc{Megalodon}-7B, \textsc{Llama}2-7B and \textsc{Llama}2-13B w.r.t processed tokens during training in Figure~\ref{fig:train_loss}.
\textsc{Megalodon}-7B obtains significantly better (lower) NLL than \textsc{Llama}2-7B under the same amount of training tokens, demonstrating better data efficiency. Moreover, \textsc{Megalodon} suffers less training spikes than the Transformer-based architecture in \textsc{Llama}2. 
Note that at the first 1/4 of the pretraining process ($<500$B tokens), the NLL of \textsc{Megalodon}-7B is slightly worse than \textsc{Llama}2-7B. We found that the main reason is that we increased the base $\theta$ of RoPE from $10,000$ in \textsc{Llama}2 to $100,000$ in \textsc{Megalodon}, which slows down model convergence at the beginning of the pretraining process. At the end, \textsc{Megalodon} reaches a training loss of 1.70, landing mid-way between \textsc{Llama}2-7B (1.75) and \textsc{Llama}2-13B (1.67). 

For computation efficiency, we conduct experiments of running \textsc{Llama}2-7B and \textsc{Megalodon}-7B using the same amount of computational resources and comparing their training speed under various context lengths. Specifically, we execute each experiment to train a model with global batch size 4M tokens distributed on 256 NVIDIA A100 GPUs (16K tokens per A100) and calculate the word/token per second (WPS) to measure the training speed. 
Figure~\ref{fig:wps} illustrates the average WPS per device of \textsc{Llama}2-7B and \textsc{Megalodon}-7B using 4K and 32K context lengths, respectively. For \textsc{Llama}2 models, we accelerate the computation of full attention with Flash-Attention V2~\citep{dao2024flashattention}. Under 4K context length, \textsc{Megalodon}-7B is slightly slower (about $6\%$) than \textsc{Llama}2-7B, due to the introduction of CEMA and Timestep Normalization. When we scale up context length to 32K, \textsc{Megalodon}-7B is significantly faster (about $32\%$) than \textsc{Llama}2-7B, demonstrating the computation efficiency of \textsc{Megalodon} for long-context pretraining. In addition, \textsc{Megalodon}-7B-32K, which utilizes chunk parallelism~(\S\ref{subsec:parallelism}), achieves about $94\%$ utilization of \textsc{Megalodon}-7B-4K.

\subsection{Short-Context Evaluation on Academic Benchmarks}

We compare \textsc{Megalodon}-7B to \textsc{Llama}2 models on standard academic benchmarks with short contexts ($< 4$K tokens), closely following the settings in \textsc{Llama}2~\citep{touvron2023llama}.
The benchmarks are grouped into the categories listed below:

\begin{itemize}[leftmargin=*]
    \item \textbf{Commonsense Reasoning} (0-shot): HellaSwag~\citep{zellers2019hellaswag}, PIQA~\citep{bisk2020piqa}, SIQA~\citep{sap2019socialiqa}, WinoGrande~\citep{sakaguchi2021winogrande}, ARC-e and -c~\citep{clark2018think}.
    \item \textbf{World Knowledge} (5-shot): NaturalQuestions (NQ, \citet{kwiatkowski2019natural}) and TriviaQA (TQA, \citet{joshi2017triviaqa}).
    \item \textbf{Reading Comprehension} (0-shot): BoolQ~\citep{clark2019boolq}.
    \item \textbf{Popular aggregated results} (5-shot): MMLU~\citep{hendrycks2020measuring}.
\end{itemize}

Table~\ref{tab:benmarks} summarizes the results of \textsc{Megalodon} and \textsc{Llama}2 on these academic benchmarks, together with other open-source base models, including MPT~\citep{mpt7b}, RWKV~\citep{peng2023rwkv}, Mamba~\citep{gu2023mamba}, Mistral~\citep{jiang2023mistral} and Gemma~\citep{team2024gemma}. Pretrained on the same 2T tokens, \textsc{Megalodon}-7B surpasses \textsc{Llama}2-7B across all the benchmarks. On some tasks, \textsc{Megalodon}-7B achieves comparable or even better performance with \textsc{Llama}2-13B. 
Note that Mistral-7B and Gemma-8B were pretrained on much larger datasets than \textsc{Megalodon}-7B, hence the results are not directly comparable.

\subsection{Long-Context Evaluation}

\paragraph{Perplexity over Long Sequences} To demonstrate the capability of \textsc{Megalodon} to make use of very long contexts to improve next-token prediction, we start by conducting the evaluation of valid perplexity on different context lengths. Concretely, we construct a validation dataset which consists of 1,920 selected books. Each of these books contains sequences with at least 2M tokens. The validation dataset is constructed by first randomly shuffling all the files and then concatenating them. 
Figure~\ref{fig:ctx-ppl} shows the perplexity (PPL) of the validation dataset in various context lengths ranging from 4K to 2M. We observe that the PPL decreases monotonically with context length, validating the effectivenss and robustness of \textsc{Megalodon} on modeling extremely long sequences.

\paragraph{Long-Context QA tasks in Scrolls} Next, we evaluate \textsc{Megalodon} on long-context open-book question answering (QA) tasks in the Scrolls dataset~\citep{shaham-etal-2022-scrolls}, including NarrativeQA~\citep{kocisky-etal-2018-narrativeqa}, Qasper~\citep{dasigi-etal-2021-dataset} and QMSum~\citep{zhong-etal-2021-qmsum}. Following \citet{xiong2023effective}, we use a simple prompt \mbox{\{\textup{CONTEXT}\} Q: \{\textup{QUESTION}\} A:} for all the tasks, and evaluate 0-shot F1-score on NarrativeQA, 2-shot F1-score on Qasper and 1-shot geometric-ROUGE\footnote{Geometric mean of ROUGE-1, 2 and L.} on QMSum.
Table~\ref{tab:scrolls} lists the results of \textsc{Megalodon}-7B, together with other open-source long-context models in the scale of 7B, namely Xgen-7B-8K~\citep{nijkamp2023xgen7b}, MPT-7B-8K~\citep{mpt7b}, YaRN-7B-128k~\citep{peng2024yarn}, \textsc{Llama}2-7B-4K~\citep{touvron2023llama} and \textsc{Llama}2-7B-32K (\textsc{Llama}2-L, \citet{xiong2023effective}).
\textsc{Megalodon}-7B obtains the best F1 on NarrativeQA, and competitive results with \textsc{Llama}2-7B Long. It should be noticed that \textsc{Llama}2-7B Long extends the context length of \textsc{Llama}2-7B from 4K to 32K by continually pretraining it on additional 500B tokens from long-context data.

\subsection{Instruction Finetuning} To evaluation the generalization capability of \textsc{Megalodon} on instruction following and alignment, We finetune the base model of \textsc{Megalodon}-7B on a proprietary instruction-alignment data under a controlled setting. We did not apply any RLHF techniques to further finetune it.
Table~\ref{tab:mtbench} summarizes the performance of chat models in 7B scale on MT-Bench\footnote{\url{https://klu.ai/glossary/mt-bench-eval}}. 
\textsc{Megalodon} exhibits superior performance on MT-Bench compared to Vicuna~\citep{vicuna2023}, and comparable performance to \textsc{Llama}2-Chat, which utilizes RLHF for further alignment finetuning. We present some outputs from instruction finetuned \textsc{Megalodon} in Appendix~\ref{appendix:demo}.

\subsection{Evaluation on Medium-Scale Benchmarks}

\paragraph{ImageNet Classification} To evaluate \textsc{Megalodon} on image classification task, we conduct experiments on the Imagenet-$1$K~\citep{deng2009imagenet} dataset, which consists of 1.28M training images and 50K validation images from 1000 classes. We mostly follow DeiT's approach of applying several data augmentation and regularization methods that facilitate the training process, and adopt most the hyperparameters from \citet{ma2023mega}.
For classification task, we replace the timestep normalization with the standard group normalization method. Top-1 accuracy on the validation set is reported in Table~\ref{tab:imagenet} to assess various models.
\textsc{Megalodon} obtains about $1.3$\% accuracy improvement over DeiT-B~\citep{touvron2021training}, and $0.8$\%. improvement over \textsc{Mega}~\citep{ma2023mega}.

\paragraph{Auto-regressive Language Modeling on PG-19} We also evaluate \textsc{Megalodon} on auto-regressive language modeling on the medium-scale PG19~\citep{raecompressive2019} datasets. We use the same vocabulary from Block-Recurrent Transformer~\citep{hutchins2022block} and adopt most of its hyper-parameters to train a \textsc{Megalodon} model with 1.3B parameters.
Table~\ref{tab:pg19} illustrate the word-level perplexity (PPL) of \textsc{Megalodon} on PG-19, together with previous state-of-the-art models, including Compressive Transformer~\citep{rae2020compressive}, Perceiver AR~\citep{hawthorne2022general}, Block-Recurrent Transformer~\citep{hutchins2022block} and \textsc{MegaByte}~\citep{yu2024megabyte}. 
\textsc{Megalodon} significantly outperforms all the baselines.

\section{Conclusion}

We have introduced \textsc{Megalodon}, an improved \textsc{Mega} architecture with multiple novel technical components, including complex exponential moving average (CEMA), the timestep normalization layer, normalized attention and pre-norm with two-hop residual configuration, to improve its capability, efficiency and scalability. Through a direct comparison with \textsc{Llama}2, 

\textsc{Megalodon} achieves impressive improvements on both training perplexity and across downstream benchmarks. Importantly, experimental results on long-context modeling demonstrate \textsc{Megalodon}'s ability to model sequences of unlimited length. Additional experiments on small/medium-scale benchmarks across different data modalities illustrate the robust improvements of \textsc{Megalodon}, which lead to a potential direction of future work to apply \textsc{Megalodon} for large-scale multi-modality pretraining.

\begin{ack}
We thank Sadhika Malladi, Zihao Ye, Dacheng Li and Rulin Shao for their helpful feedback and discussion during this work.
\end{ack}

\end{document}
